\chapter{Research Methodology}
\label{chap:methodology}

\section{Research Design and Status}
\label{sec:method-status}

A controlled longitudinal experiment is proposed to compare delayed browser profiling with evidence retained during the original visit. Neither the finite feasibility pilot, FP-0.1, nor the main study has been run. The procedures and numerical settings in this chapter are proposed protocol choices. Actual host capacity, software builds, capture fields and classifier mappings remain to be verified before collection.

The main unit of evaluation will be an \emph{original regression episode}. Each affected scheduled original visit will receive an episode identifier and will be linked to any diagnostic follow-up. Visits affected by one continuing fault occurrence will retain that shared occurrence identifier. They will not be treated as independent sites. Each strategy will receive a separate execution within a matched block; matching will align the assigned fixture, scenario and repetition, but cannot make nondeterministic executions identical \parencite{GJMALCNF,E8MGSAC8,B9FIJAKE}.

The primary comparison will be retained original-event evidence (C) minus detailed follow-up profiling (B) for intermittent faults under common resource allowances and a diagnosis deadline. Objectives O1--O5 will be addressed through the measurement contract, labelled dataset, strategy comparison, outcome analysis and assessment of reliability conditions. The pilot will assess feasibility. The main-study sample, target population, weighting, practical effect criterion, uncertainty method and stopping rule will be fixed separately after the pilot.

\section{O1: Measurement and Visit Contract}
\label{sec:o1-contract}

\subsection{Environment and Visit Boundaries}

Chromium, automation, profiler and metric-library versions will be fixed and recorded with checksums. The host, operating system (OS), background load, viewport, CPU and network settings, browser process and profile, cache and session state, and service-worker policy will be recorded. Intended and observed conditions will be stored separately. A fresh browser context isolates browser storage, but does not establish that server, host, Domain Name System or all cache state has been reset \parencite{SSS4NSXC,UKNA4QJZ,92RRK4LT}.

FP-0.1 is limited to two owned localhost fixtures, one headless Chromium build, one automation stack and one host. A desktop viewport of 1365 by 768 pixels, device scale factor 1, no network throttling and CPU rate 1 are specified, except for the nuisance condition. Service workers, authentication, human participants, mobile devices, live websites and multiple browsers are outside the pilot. These limits apply to feasibility work, not to the complete engineering deliverables.

The visit clock will start at attempted browser-process launch. Collection will start before fixture code executes. Each visit will be observed for 15 seconds from navigation start, with a safe button input attempted at 5 seconds. An unavailable input will be recorded as a failure. Collectors will be stopped at 15 seconds and flushed while the page remains available. Startup, observation, export and cleanup must fit within the 30-second visit limit. Original and follow-up visits will have separate identifiers, local time origins and explicitly aligned clocks.

\subsection{Metrics and Capture Definitions}

Largest Contentful Paint (LCP), measured in milliseconds, will be the primary pilot symptom. Its candidate, input cutoff and finalization rule will be fixed to the selected implementation. Client-observed navigation and resource timings and supported attribution will also be recorded. The navigation TTFB start point, timestamp formula, redirects and cache scope will be specified before collection; subresource intervals will remain separate. Interaction to Next Paint (INP) and Cumulative Layout Shift (CLS) will be exploratory checks of availability and lifecycle handling. Missing or unsupported values will receive reason codes, not zeros. Total Blocking Time (TBT), composite scores, memory-growth diagnosis and full metric validation are outside the pilot endpoint. TBT will not be treated as a substitute for INP, and historical Time to Interactive (TTI) will not be restored as an endpoint \parencite{QF52VUCU,QHUCFVBF,9T8ACIJ8,N33VTKAN,JJWRMF5Z,NXHCE5SV}.

Each capture contract will identify the collector, event types and fields, frame and process scope, start and stop times, sampling or thresholds, buffer limits, loss handling, serialization and costs. The candidate rich contract will combine Chrome DevTools Protocol (CDP) Network metadata with main-thread, paint, layout and script trace events. A sampled CPU profile will be included only if the selected build supports the chosen configuration. Response bodies, video and heap snapshots will be excluded. The exact categories and fields will be qualified before the contract is treated as implemented. A warning-free trace will not be assumed complete \parencite{REP5RHB6,63QKVL86,UMVA36NP,VMWYVJKL}.

\begin{table}[htbp]
\centering
\caption{Minimum measurement-contract fields. Interface documentation identifies candidate evidence; runtime qualification will establish the implemented contract \parencite{REP5RHB6,N33VTKAN,63QKVL86}.}
\label{tab:measurement-contract}
\begin{tabular}{p{0.22\linewidth}p{0.69\linewidth}}
\toprule
Record & Required fields \\
\midrule
Environment and journey & Builds and hashes; host, viewport and profile; intended and actual state; fixture and navigation; input, observation and finalization times. \\
Evidence and missingness & Visit identifier and role; collector, categories and scope; drop and eviction intervals; export and hash; unsupported, absent, failed, timed-out or lost reason codes. \\
Resource cost & Elapsed time, CPU time, resident memory and bytes by stage; attempts, follow-ups, failures and remaining allowance. \\
\bottomrule
\end{tabular}
\end{table}

\subsection{Conformance and Instrumentation Effects}

Conformance checks will link known events to the expected visit, timestamp, resource or script. Metric definitions, callback finalization, export and loss signals will be checked on simple fixtures. Comparisons with Chrome DevTools or Lighthouse will assess consistency where their definitions and conditions match. Since these tools share Chromium evidence, agreement will not supply an independent fault oracle \parencite{AAW3TPVK,4MGN3SYQ,VMWYVJKL,9WMJW57H}.

A minimally observed contract, M, will contain a minimal LCP observer, harness timing and the common hidden verification probe. Its measurement cost is nonzero. Separate M, A, C and rich executions will be compared in randomized calibration blocks to assess timing changes, failures and altered fault activity. The proposed calibration has six blocks per fixture: four stable, one R and one J. The stable blocks will provide a small-sample LCP perturbation summary; the fault blocks will check activity under every contract. The resulting measurement contract, conformance matrix and state-reset procedure will be versioned. Unsupported endpoints will be excluded from the main contract unless their failures are resolved.

\section{O2: Longitudinal Dataset and Fault Oracle}
\label{sec:o2-dataset}

\subsection{Fixtures and Fault Scenarios}

Fixture F1 will be an image-led page with a designated LCP candidate. F2 will use a small script to construct its designated LCP content and will have a different Document Object Model (DOM) and loading structure. Both will provide a safe button input and support both fault classes. External dependencies, cross-origin frames and background polling will be excluded. URLs, element identifiers and content structure will remain unchanged across fault conditions. Policy-visible visit identifiers will be opaque and will not encode fault names or schedules.

Fault R will delay a designated critical response by an additional 500 ms. Its actual execution will be checked using evaluator-only request identifiers and start/end logs. Fault J will introduce approximately 500 ms of main-thread work before the designated LCP content can paint. A common evaluator probe will record its entry, exit and elapsed work; the probe's cost will be measured. Probe records will be sent through a separate evaluator channel that policy processes cannot read. Policy inputs and saved evidence will be audited for accidental fault-marker leakage.

Both faults are intended to produce similar loading symptoms through different mechanisms. Qualification will establish whether each fault executes and produces the intended effect under minimally observed conditions. Performance mutation testing provides a precedent for controlled performance changes, but its native-program experiments do not validate these browser faults or supply transferable repetition counts \parencite{WF5NM72E}. A configuration flag or a rich trace alone will not be accepted as proof of fault activity.

\begin{table}[htbp]
\centering
\caption{Proposed FP-0.1 scenarios. Follow-up activity will be checked independently on each visit.}
\label{tab:pilot-scenarios}
\begin{tabular}{p{0.29\linewidth}p{0.28\linewidth}p{0.32\linewidth}}
\toprule
Scenario & Original visit & B follow-up, if alerted \\
\midrule
Stable & No fault; CPU rate 1 & No fault; CPU rate 1 \\
Nuisance & No software fault; CPU rate 2 & Same CPU setting; no fault \\
Persistent R / J & R / J & Same assigned class \\
Transient R / J & R / J & No fault \\
Replaced R / J & R / J & J / R respectively \\
\bottomrule
\end{tabular}
\end{table}

The paired rows in Table~\ref{tab:pilot-scenarios} represent eight scenarios in total. Persistence will be controlled by visit role, with a separately reset occurrence for each arm. This tests known conditions; it does not estimate natural fault lifetimes. Repeated slowness on a follow-up will not establish that the original cause recurred.

\subsection{Stable History, Alert and Diagnosis Rules}

For each fixture, the first 20 stable A visits will train the baseline and the next 20 will be used for validation. The proposed alert threshold is the median stable LCP plus the greater of 100 ms and three times the unscaled median absolute deviation (MAD). This operational threshold is neither a significance test nor a field performance threshold. It will be locked after qualification and will not be updated during the comparison.

A small rule-based classifier will be fixed after qualification. Its candidate discriminators are a critical-resource phase at least 200 ms above its stable median and at least 200 ms of attributable pre-paint main-thread work. Availability, overlap and ambiguity will be checked, and the final field-to-class mapping will be saved before comparison. Supported attribution available to A will be retained even when it already distinguishes the faults. If both explanations are supported, neither is supported or the evidence is incomplete, the classifier will abstain.

All arms will issue cause claims only after the common alert. The vocabulary will be R, J or abstain; controls will also allow ``no supported target fault.'' For B, supported original evidence will take priority. A class supported only by the follow-up will be marked as an inference about the original cause from a later visit. Conflicting original and follow-up evidence will be flagged rather than overwritten.

\subsection{Independent Labels, Eligibility and Provenance}

An evaluator-only intervention manifest and per-visit activity and applicability checks will define the labels. An eligible original opportunity will be a scheduled visit in a prequalified fault case whose fault execution and declared journey/endpoint applicability are independently confirmed. Eligibility will not depend on an alert, a slow observed value, metric availability or successful capture. A fault-induced timeout, missing policy metric or missed alert will not remove an otherwise eligible opportunity.

Cases with unresolved activity or cause will remain in the assignment ledger with an unresolved-label status, outside known-cause scoring. The same eligibility rule will be applied to all arms. Predictions will be frozen before hidden labels are joined for evaluation. Intended interventions, verified execution and policy diagnoses will remain separate records. This preserves the distinction between what was scheduled, what occurred and what the strategy inferred.

\begin{table}[htbp]
\centering
\caption{Linked dataset records. Explicit units, state and missingness support reproducibility; the fault oracle is a study-specific design \parencite{RB93JMBB,B9FIJAKE,WF5NM72E}.}
\label{tab:dataset-schema}
\begin{tabular}{p{0.22\linewidth}p{0.69\linewidth}}
\toprule
Record & Minimum contents \\
\midrule
Site and configuration & Template, deployment, journey, versions, host and profile; intended and observed state. \\
Visit and episode & Visit, parent episode, block, occurrence and session identifiers; role and timings; class, onset and duration; activity, persistence, applicability and label status. \\
Evidence and diagnosis & Raw file references and hashes; capture contract, window and loss; prediction or abstention, diagnosis time and evidence available to the policy. \\
Cost and attempt & Resource dimensions, stage, outcome and failure codes, remaining allowance, follow-ups and evaluator-only cost. \\
\bottomrule
\end{tabular}
\end{table}

An all-attempt record will separate assigned, attempted, verified eligible, known ineligible, unresolved-label and globally unattempted opportunities. Raw evidence, derived metrics and labels will be stored separately. For the main study, historical training, tuning and later test periods will be separated, and distinct fixture templates or site instances will be reserved before tuning. Pilot and repair observations will be excluded from the confirmatory dataset.

\section{O3: Capture Strategies and Allocation}
\label{sec:o3-strategies}

\subsection{Evidence Available to Each Strategy}

All strategies will share the lightweight measurements, baseline, alert rule, fault vocabulary, scheduled original opportunities and diagnosis deadline. Each will receive its own execution. Arm order will be randomized within scenario blocks, the specified state will be reset and only one browser workload will run at a time. Blocking controls assigned factors but does not remove every source of carryover \parencite{FBTV2ICM,GJMALCNF}.

Figure~\ref{fig:evidence-chain} shows how evidence and labels will be separated. A locked diagnosis will be joined to the verified original label only at evaluation.

\begin{figure}[htbp]
\centering
% Original study-design illustration, not observed data or a copied architecture.
\begin{tikzpicture}[
  box/.style={draw,rounded corners=2pt,align=center,text width=6cm,
              minimum height=1.0cm,inner sep=5pt,font=\small},
  side/.style={box,text width=4.4cm},
  link/.style={-{Latex[length=2mm]},thick}]
\node[box] (original) {Assigned original visit\\Opaque visit and episode identifiers};
\node[box,below=7mm of original] (policy) {Policy evidence: legitimate original attribution;\\retained entries, scheduled trace or rich follow-up};
\node[box,below=7mm of policy] (prediction) {Lock an original-cause claim or abstention\\Within the common deadline and resource allowance};
\node[box,below=7mm of prediction] (score) {Join labels only after prediction lock\\Score every verified eligible original episode};
\node[side,right=1cm of prediction] (oracle) {Evaluator-only checks\\Verify original and follow-up fault activity separately\\Hidden labels never enter policy rules};
\node[side,right=1cm of score,yshift=-5mm] (ledger) {Attempt and cost ledgers\\Retain failures, loss, uncertainty and unattempted assignments};
\draw[link] (original)--(policy);
\draw[link] (policy)--(prediction);
\draw[link] (prediction)--(score);
\draw[link] (original.east)-|(oracle.north);
\draw[link,dashed] (oracle.south west)--(score.north east);
\draw[link] (ledger.west)--(score.east);
\end{tikzpicture}

\caption{Proposed evidence and hidden-label flow for scoring the cause of the original episode.}
\label{fig:evidence-chain}
\end{figure}

\begin{table}[htbp]
\centering
\setlength{\tabcolsep}{3pt}
\caption{Proposed capture strategies. Adaptive tracing and retained original evidence are established mechanisms; their effectiveness and costs in this comparison remain to be measured \parencite{A3SS4UWX,3ED3FDXH,TIIDVPV8,BUPGR8ZJ,N33VTKAN}.}
\label{tab:arms}
\begin{tabular}{p{0.12\linewidth}p{0.57\linewidth}p{0.25\linewidth}}
\toprule
Arm & Evidence contract & Role \\
\midrule
A & Metrics and all supported original-visit attribution, serialized before fields disappear. & Practical baseline. \\
B & A on the original visit; one rich follow-up after an alert, with A evidence retained. & Delayed diagnosis. \\
C & A plus an early record of supported Resource Timing, long-animation-frame and long-task entries; diagnostic access after alert. & Bounded original evidence. \\
D & A on every original; rich capture on every fourth assigned original per fixture, with randomized starting phase. & Periodic alternative. \\
E & A and rich capture on every original, under a separate reference allowance. & Coverage/cost reference. \\
\bottomrule
\end{tabular}
\end{table}

C will retain at most 1,024 entries or 4 MiB of serialized data, whichever limit is reached first. Evictions and drops will be recorded with affected time ranges. At 15 seconds, its record will be stopped and frozen before alert finalization, then exported within the visit deadline. Every C original, including an unalerted one, will be exported for conformance auditing to a separate, access-logged store inaccessible to the diagnosis worker. Every export will be charged. A diagnostic copy will be released only after the common alert; the classifier will have no early access through the audit store. Missing required earlier evidence will make the record incomplete. This limited entry record is not a full browser trace.

D's phase will be selected independently of the hidden fault schedule, and no additional visit will be introduced. E will neither provide labels nor share evidence with other arms. Its richer observation can itself perturb the workload or miss evidence, so it is not an oracle or guaranteed upper bound. Hindsight's bounded retention lifecycle motivates explicit horizon and loss checks, but its distributed-system capacities do not prescribe browser buffer sizes \parencite{3ED3FDXH,REP5RHB6}.

\subsection{Finite Allocation and Schedule}

FP-0.1 specifies 64 matched blocks: two fixtures, eight scenarios and four repetitions. Each block will contain five separate originals, one per arm, giving 320 comparison originals. Each arm will receive 48 assigned fault opportunities, eight stable controls and eight nuisance controls. Transient and replaced-cause cases will each contribute 16 originals per arm. These are assigned counts, not guaranteed completed or verified observations.

\begin{table}[htbp]
\centering
\caption{Maximum planned browser starts in FP-0.1. All counts are feasibility allocations, not a power calculation.}
\label{tab:pilot-allocation}
\begin{tabular}{p{0.29\linewidth}p{0.52\linewidth}r}
\toprule
Stage & Allocation & Starts \\
\midrule
Warm-up & Two visits per fixture in each of three sessions & 12 \\
Stable history & Two fixtures, 40 A visits each & 80 \\
Qualification & Two fixtures, four conditions, M/C/rich contracts, five visits; four extra C checks per fixture & 128 \\
Perturbation calibration & Two fixtures, six blocks, four contracts & 48 \\
Comparison originals & Two fixtures, eight scenarios, four repetitions, five arms & 320 \\
B follow-ups & At most one per B original, including control alerts & 64 \\
\midrule
Planned maximum & Preparation and comparison combined & 652 \\
Repair reserve & One pre-comparison repair cycle only & 48 \\
Absolute cap & No additional or replacement visits & 700 \\
\bottomrule
\end{tabular}
\end{table}

Warm-up records will be retained but excluded from evaluation. Preparation and calibration will occupy one session of at most three active hours. Two comparison sessions on different days will each have a 2.5-hour limit. Repetitions 1--2 will be allocated to the first comparison session and 3--4 to the second. Fixture/scenario blocks and arm order within each block will be shuffled. Seeds and schedule hashes will be saved before comparison, and the hidden fault schedule will be withheld from policies.

B will be assigned a target launch delay of 1 or 15 seconds after the original 15-second observation window, with each delay used twice across the four repetitions in a fixture/scenario cell. Actual delay to follow-up navigation, including startup and queueing, will be recorded. The common diagnosis deadline will be 90 seconds after the original observation window ends. Scheduling feasibility, especially the 1-second launch target alongside export and cleanup, will be checked before collection. Missed targets and late diagnoses will remain outcomes; further diagnostic visits will not be used to repair them.

Four repetitions allow both delays to be represented across two sessions; they do not establish population precision. The two fixtures remain two templates despite the larger number of visits. Published five-run Lighthouse recommendations and native-benchmark repetition tables will not be transferred as sample-size rules \parencite{DJI5ULJD,GJMALCNF}.

\subsection{Resource Accounting and Stops}

The campaign will stop at the first of 700 starts, eight active hours, seven consecutive calendar days from the first warm-up or 8 GiB of artifacts. Active time includes startup, waits, resets, verification, capture and export. Gaps between sessions count towards the calendar limit. MiB and GiB denote binary storage units. Resident set size (RSS) will refer to memory resident in the browser and collector process tree.

\begin{table}[htbp]
\centering
\caption{Proposed resource limits and responses. A--D have equal allowance vectors; E has a separate reference allowance.}
\label{tab:resource-limits}
\begin{tabular}{p{0.26\linewidth}p{0.65\linewidth}}
\toprule
Scope & Limit and response \\
\midrule
Each visit & 30 seconds elapsed, 30 CPU-seconds and 2 GiB peak process-tree RSS; a breach will terminate the visit and retain its reason. \\
Rich export & 32 MiB per visit; truncation and loss will be marked. \\
C record & 1,024 entries or 4 MiB serialized data; runtime memory will be measured separately. Serialized size is not heap usage. \\
Each A--D policy & Three charged hours, 120 CPU-minutes and 2 GiB cumulative exported evidence, with the common visit limits. \\
E reference & Three charged hours, 120 CPU-minutes and 3 GiB exported evidence, with the common visit limits. \\
Diagnosis and retries & Common 90-second deadline; at most one B follow-up; no comparison retries or visits until recurrence. \\
\bottomrule
\end{tabular}
\end{table}

All original scans, retained or dropped records, serialization, exports, waits, confirmations, failed attempts and follow-ups will be charged. CPU time, elapsed time, memory and bytes will remain separate quantities. The full shared stable-history cost will be charged to each arm's standalone policy ledger, but its physical visits will be counted only once in the campaign ledger. Qualification and calibration will be reported as development costs outside policy allowances and inside campaign caps. Evaluator-only verification will remain hidden and will be costed separately.

After a policy exhausts its allowance, it will stop obtaining evidence. Remaining scheduled originals will still be executed with the minimally observed evaluator harness while campaign limits permit. Verified eligible episodes will remain in that policy's denominator as no correct in-time diagnosis. In contrast, a global cap or shared integrity stop will prevent remaining launches. Those cases will be marked unattempted without invented labels, and the campaign will be reported as incomplete. Evidence will not be deleted during collection to make a breached storage cap appear satisfied.

\section{O4: Outcomes, Denominators and Analysis}
\label{sec:o4-analysis}

\subsection{Primary Comparison and Missing Labels}

A correct attribution will name the independently verified fault class of the original episode within the common deadline and policy allowance. A follow-up-only explanation will be scored against the original label, even if the follow-up contains a different fault.

For arm $s$, let $E_s$ contain its independently verified eligible original regression episodes and let $N_s=|E_s|$. Let $c_s$ count episodes in $E_s$ with a timely, within-budget, locked, non-abstaining prediction of the correct original class. The attribution rate and C-minus-B contrast will be
\begin{equation}
\widehat{p}_s=\frac{c_s}{N_s},\qquad
\widehat{\Delta}_{C-B}=\widehat{p}_C-\widehat{p}_B.
\label{eq:primary-attribution}
\end{equation}
The principal contrast will use a prospectively defined intermittent subset. Pilot transient and replacement cases will be shown separately. A rate will be undefined when $N_s=0$. Missed alerts, failed or late follow-ups, lost evidence, policy exhaustion and abstentions will remain in $N_s$ as no correct diagnosis. Repeated alerts for one original episode will be deduplicated.

The five executions within a matched block can have different verified eligibility or unresolved-label patterns. Matching therefore does not imply equal observed denominators. The main-study analysis freeze will define the target population and site/scenario weighting, the treatment of unequal verified arm sets, and the sensitivity analysis for unresolved labels. No pairing or weighting method is assumed solely from the matched allocation. Pilot counts will be shown by arm and cell so that differential verification cannot disappear in a pooled rate.

For the pilot, both the 48 assigned fault opportunities per arm and the verified eligible denominator will be reported. An assignment-based sensitivity range will use known correct cases divided by 48 as its lower bound. Only attempted cases with unresolved eligibility or cause will be added to the numerator for the upper bound. Corresponding worst/best combinations will bound the descriptive C--B difference; transient and replacement subsets will use their own assigned denominators. These missing-label bounds are not confidence intervals and cannot turn a globally incomplete campaign into a completed comparison. Globally unattempted cases will remain separately identified.

\subsection{Secondary Outcomes and Interpretation}

Detection coverage will be the fraction of verified eligible episodes alerted. Diagnostic coverage will be the fraction receiving a timely, within-budget, non-abstaining cause claim. Conditional diagnostic precision will be correct original-cause claims divided by all such claims on verified eligible episodes; it will be undefined if no claims are made. Wrong original-cause claims will be reported over all verified eligible episodes with a fault-class confusion matrix. Coverage will accompany precision so that selective abstention cannot conceal missed episodes.

Stable and nuisance controls will be reported separately, with false-alert and unsupported cause-claim rates per observation and per unit of monitoring time. Abstention, timeout, failed follow-up, unsupported output, buffer loss, unavailable evidence and exhausted allowance will be recorded with overlapping reason codes where necessary. Assigned and actual follow-up delay, independently checked recurrence and diagnosis time will be reported. Unresolved episodes will remain visible alongside the timing summaries of successful diagnoses.

Costs will be reported in total, per scheduled original and per correctly attributed episode. Cost per correct attribution will be undefined when no episode is correctly attributed. Counts and rates will be broken down by fixture, scenario, arm and session. Pilot summaries will include medians, ranges, variation across blocks and sessions, and ordering or drift plots. No efficacy significance test or two-fixture cluster bootstrap will be presented as population uncertainty.

The main study will require a practical effect or precision target, justified independent-unit coverage, a sample-size rationale, an appraised interval or test, multiplicity handling, missing-label analysis and a stopping rule. Pilot variation and costs will inform a range of sample-size scenarios rather than a calculation based on one favorable effect. Nested benchmarking and performance-fault studies motivate this care, but do not supply the proposed study's sample size or causal oracle \parencite{GJMALCNF,E8MGSAC8,WF5NM72E}. An inconclusive comparison will not establish equivalence.

\section{O5: Reliability Boundaries and Validity}
\label{sec:o5-validity}

Fault persistence will be the principal factor. Attribution and cost will be compared across persistent, transient and replaced-cause conditions. Follow-up delay and controlled browser-state changes will be secondary factors where the budget and later protocol support them. A condition-by-strategy summary will distinguish adequate attribution, evidence loss and inconclusive results using a practical criterion fixed before main-study outcomes are inspected.

The pilot will use a fresh browser process and empty profile, cleared fixture/controller state and disabled service workers. Host load and previous-run state will be recorded; complete host-cache clearing will not be assumed. A later browser-state experiment would require one-factor contrasts where feasible. Uncontrolled changes in several conditions would not establish an isolated cache effect \parencite{HS3PKWFL,ISZ9TU75,SSS4NSXC,U2NAEPYI}.

Internal-validity risks include label leakage, incorrect visit linkage, missing entries, measurement perturbation, state carryover and tuning on evaluation data. Separate evaluator access, opaque identifiers, independent activity checks, randomized order, versioned manifests and complete attempt records will address these risks. Repeated visits and related fault occurrences will retain their dependence identifiers. The small number of fixture templates limits external validity; localhost conditions cannot establish production fault prevalence or real-user accuracy \parencite{P57ARGS6,RB93JMBB,B9FIJAKE}. Any later authorized live-site sample will assess operational robustness, not supply verified fault truth from unexplained slowdowns.

\subsection{Data Ethics, Retention and Reproducibility}

The pilot will use owned localhost fixtures, fixed safe inputs and no human participants or real-user traces. Credentials and unnecessary page or user content will not be stored. Policy evidence and evaluator labels will use separate access paths, and an append-only manifest will preserve attempts, hashes, loss codes and deviations. Access logging will support the audit of C's retained evidence and hidden labels.

Within the 8-GiB cap, raw pilot artifacts will be retained for 90 days after pilot-report sign-off. The de-identified ledger, configurations, derived tables and hashes will be retained through project assessment. Actual storage and deletion actions will be recorded when due. Any later live-site work will require its own authorization, access conditions and data-minimization rules.

\section{Pilot Progression and Main-Study Separation}
\label{sec:pilot-main}

Before collection, the execution manifest will be completed with the host, OS, CPU, RAM, browser and automation builds, package checksums, fixture commits, capture categories, clock alignment, diagnosis mapping, seeds, storage location, start date and accountable researcher. Available host capacity and enforcement of the caps will be checked. Qualification will include four additional C checks per fixture: buffer overflow, export failure, a trigger delayed to 20 seconds and absence of the qualifying input. Expected failures in these checks will be separated from unexpected failures in ordinary visits.

The proposed progression rules are engineering decisions, not statistical acceptance tests. For each fixture and class under M, all five qualification injections must execute, at least four must exceed the alert threshold, and median LCP must increase by at least 200 ms relative to stable M visits. Fault activity under C and rich capture will also be checked. C must supply correctly scoped, exportable resource and main-thread discriminators in every ordinary qualification case; its frozen rules must distinguish each class in at least four of five cases per fixture. Rich-capture discriminator mapping will be checked separately.

No required-record loss is permitted for ordinary C qualification visits, and deliberately induced loss must be detected. Stable median LCP increases for both A and C relative to M must be no greater than the larger of 25 ms and 10\% of the M median for each fixture, without altered fault activity. Increases above this level but no greater than the larger of 60 ms and 25\%, or unstable estimates, will require amendment; larger changes or altered mechanisms will stop this protocol version. Rich-reference perturbation will be reported separately.

Operational progression after comparison requires every design cell to have been attempted, at least 58 of 64 originals per arm to have completed with valid exported evidence, and at least 46 of 48 assigned fault labels to have been verified. Counts of at least 48 of 64 and 39 of 48 respectively may support amendment if a specific feasible correction exists. Lower counts, an unattempted cell or an unworkable resource requirement will stop this version. Every known-cause score must have valid independent verification and linkage. Label leakage, wrong visit linkage or unreliable shared records will cause an immediate integrity stop.

Green progression means that a main-study protocol may be prepared; it does not establish efficacy or authorize automatic main collection. Overall amber requires a specific prospective amendment. This differs from the single repair cycle allowed \emph{before} comparison: at most 48 extra starts within all campaign caps, with a documented failure and correction. The only permitted severity increase is from 500 to 1,000 ms if qualification fails to produce the intended symptom. A new fault class, fixture, detector, capture mechanism or larger budget requires a new protocol version. Comparison outcomes will not be used to tune that same comparison.

A signalled arm-specific timeout, overflow or missing record will remain an outcome; other scheduled work may continue while shared integrity and caps permit. If A already diagnoses the selected causes, that result will be retained. If C cannot provide useful original evidence, its scope or the research comparison will be revised prospectively. C outperforming B is not a progression requirement.

The pilot readout is limited to six analyst-hours within two calendar days of the collection stop. It will contain run and cost ledgers, fixture and capture manifests, conformance and loss checks, per-cell outcome tables, perturbation summaries, deviations and a reasoned stop, amend or prepare-main-study decision. An unfinished readout will be recorded as incomplete without extending collection. Main-study collection will depend on a separate protocol and analysis freeze, and all pilot and repair data will remain outside the confirmatory dataset.
